\subsection{Bài Toán Data Fitting}


\subsubsection{Phát Biểu Bài Toán Tổng Quát}

\begin{tcolorbox}[mycodebox, colback=blue!3!white, colframe=blue!40!black,
    title={\textbf{Định nghĩa 1.1} - Bài Toán Data Fitting}]
Cho tập dữ liệu $\mathcal{D} = \{(\mathbf{x}_i, y_i)\}_{i=1}^n$ với
$\mathbf{x}_i \in \mathbb{R}^p$, $y_i \in \mathbb{R}$.
Bài toán Data Fitting là tìm hàm $f : \mathbb{R}^p \to \mathbb{R}$ trong một
lớp hàm cho trước sao cho $f$ xấp xỉ tốt nhất ánh xạ từ $\mathbf{x}_i$ đến
$y_i$ theo một tiêu chí đã định.
\end{tcolorbox}

Trong mô hình \textbf{hồi quy tuyến tính}, ta giả thiết quan hệ tuyến tính
giữa các biến độc lập và biến phụ thuộc:
\begin{equation}
    y_i = \beta_0 + \beta_1 x_{i1} + \beta_2 x_{i2} + \cdots + \beta_p x_{ip}
          + \epsilon_i
        = \mathbf{x}_i^T \boldsymbol{\beta} + \epsilon_i,
    \label{eq:linear_model}
\end{equation}
với $\boldsymbol{\beta} = (\beta_0, \beta_1, \ldots, \beta_p)^T \in \mathbb{R}^{p+1}$
là vector tham số cần ước lượng và $\epsilon_i$ là nhiễu ngẫu nhiên.

Viết dưới dạng ma trận, với
$\mathbf{X} \in \mathbb{R}^{n \times (p+1)}$ là \emph{ma trận design} (cột đầu
toàn 1 đại diện cho intercept):
\begin{equation}
    \mathbf{y} = \mathbf{X}\boldsymbol{\beta} + \boldsymbol{\varepsilon},
    \label{eq:matrix_form}
\end{equation}
trong đó $\mathbf{y} \in \mathbb{R}^n$ là vector quan sát và
$\boldsymbol{\varepsilon} \in \mathbb{R}^n$ là vector nhiễu.

\subsubsection{Các Giả Thiết Gauss-Markov}

Để đảm bảo OLS có các tính chất tốt, ta cần các giả thiết Gauss-Markov
(GM1-GM5):

\begin{enumerate}[label=\textbf{GM\arabic*.}]
    \item \textbf{Tuyến tính:} $\mathbf{y} = \mathbf{X}\boldsymbol{\beta} +
          \boldsymbol{\varepsilon}$.
    \item \textbf{Không hoàn hảo đa cộng tuyến:} $\rank(\mathbf{X}) = p+1$
          (các cột của $\mathbf{X}$ độc lập tuyến tính, đảm bảo
          $\mathbf{X}^T\mathbf{X}$ khả nghịch).
    \item \textbf{Ngoại sinh:} $E[\boldsymbol{\varepsilon} \mid \mathbf{X}] =
          \mathbf{0}$, tức $E[\epsilon_i \mid \mathbf{x}_i] = 0$ với mọi $i$.
    \item \textbf{Đồng phương sai (Homoscedasticity):}
          $\mathrm{Var}(\boldsymbol{\varepsilon} \mid \mathbf{X}) =
          \sigma^2 \mathbf{I}_n$, tức các sai số có cùng phương sai $\sigma^2$
          và không tương quan với nhau.
    \item \textbf{Phân phối chuẩn (tùy chọn):}
          $\boldsymbol{\varepsilon} \mid \mathbf{X} \sim
          \mathcal{N}(\mathbf{0}, \sigma^2 \mathbf{I}_n)$
          (cần cho kiểm định giả thuyết).
\end{enumerate}


